% Browser-Native Spiking Neural Networks in Contracted Physics
% NeurIPS 2027 Workshop on Neuromorphic Computing / Embodied AI
% Paper ID: P1-SNN (SNN compute path under contract in Program 1: Trust-by-Construction)
%
% Build: pdflatex paper-2-snn-neurips.tex
%        bibtex paper-2-snn-neurips
%        pdflatex paper-2-snn-neurips.tex (x2)

\documentclass[10pt]{article}

% --- NeurIPS-style formatting ---
% Guarded load: matches the USENIX .sty pattern used by papers 1 and 4
% (see paper-audit-matrix.md venue-fit note, 2026-04-19). When the vendor
% neurips_<year>.sty is present on submission, this picks it up cleanly.
% When absent (local compile without the style package), we fall back to a
% reasonable default (1in margins, Times). Bundle-time step on submission:
% ship neurips_2027.sty alongside the .tex.
\IfFileExists{neurips_2027.sty}{%
  \usepackage[final]{neurips_2027}%
}{%
  \typeout{Skipping neurips_2027: vendor sty absent; local compile uses default article class.}%
  \usepackage[margin=1in]{geometry}%
  \usepackage{times}%
  \setlength{\parskip}{0.5em}%
}

\usepackage[utf8]{inputenc}
\usepackage[T1]{fontenc}
\usepackage{hyperref}
\usepackage{url}
\usepackage{booktabs}
\usepackage{amsfonts}
\usepackage{amsmath}
\usepackage{amssymb}
\usepackage{amsthm}
\usepackage{enumitem}  % for [leftmargin=*,topsep=2pt] enumerate options
\newtheorem{theorem}{Theorem}
\newtheorem{lemma}{Lemma}
\newtheorem{definition}{Definition}
\newtheorem{property}{Property}
\newtheorem{corollary}{Corollary}
\usepackage{nicefrac}
\usepackage{microtype}
\usepackage{graphicx}
\usepackage{xcolor}
\usepackage{algorithm}
\usepackage{algorithmic}
\usepackage{listings}
\usepackage{subcaption}
\usepackage{tikz}
\usetikzlibrary{arrows.meta, positioning, shapes.geometric, calc}

% --- Custom commands ---
\newcommand{\todo}[1]{\textcolor{red}{\textbf{[TODO: #1]}}}
\newcommand{\snnwebgpu}{\textsc{SNN-WebGPU}}
\newcommand{\holoscript}{\textsc{HoloScript}}
\newcommand{\cael}{\textsc{CAEL}}
\newcommand{\relu}{\operatorname{ReLU}}
\newcommand{\tropplus}{\oplus}
\newcommand{\tropmult}{\otimes}
\DeclareMathOperator*{\argmax}{arg\,max}

% Lotus petal data contract — \measuredFrom{<artifact>} marks values that come
% from a runnable artifact in the repo. Renders nothing in the PDF; grep target
% for scripts/build-lotus.mjs (research/2026-04-26_lotus-petal-data-contract.md §3+§4).
\providecommand{\measuredFrom}[1]{}
\providecommand{\measuredAt}[1]{}

% --- Listings style for WGSL ---
\lstdefinelanguage{WGSL}{
  keywords={fn, let, var, if, else, for, return, struct, const, true, false},
  morekeywords={f32, u32, i32, vec3, array, compute, workgroup_size, builtin, group, binding, uniform, storage, read, read_write},
  sensitive=true,
  morecomment=[l]{//},
  morestring=[b]",
}
\lstset{
  language=WGSL,
  basicstyle=\ttfamily\scriptsize,
  keywordstyle=\color{blue}\bfseries,
  commentstyle=\color{gray}\itshape,
  numbers=left,
  numberstyle=\tiny\color{gray},
  breaklines=true,
  frame=single,
  captionpos=b,
}

% ============================================================================
\title{Browser-Native Spiking Neural Networks in Contracted Physics:\\
       WebGPU Compute Shaders for Embodied Neuromorphic Cognition}

\author{
  Joseph Krzywoszyja \\
  Infinitus Systems \\
  \texttt{josep@holoscript.net} \\
}

\begin{document}

\maketitle

% ============================================================================
\begin{abstract}
We present \snnwebgpu{} as the first executable Tier~1 of an \emph{inverted
runtime stack}: neuromorphic computation is the primary executor, not an
accelerator bolted onto a CPU-centric engine, and every spike decision is backed
by a cryptographic determinism receipt that makes NN-primary runtime stacks
architecturally honest.
The system implements Leaky Integrate-and-Fire (LIF) neurons, Spike-Timing-Dependent
Plasticity (STDP) and multi-layer SNN architectures entirely in WGSL compute
shaders dispatched from a standard web browser via the WebGPU API, requiring
zero installation and no server-side computation.
We pair a small, fast SNN proposer with a CPU-side verifier
(SimulationContract) generalizing the 2026 speculative-decoding pattern from
token-stream to operation-stream; on neuromorphic-class hardware this inversion
delivers single-digit milliwatt steady-state power---three orders of magnitude
below GPU-accelerated LLM inference.
We establish a formal bridge between SNN rate coding and classical ReLU activations
via tropical algebra ($\max(0,x) = x \tropplus 0$ in the max-plus semiring),
enabling principled interoperability between spiking and non-spiking networks.
Additionally, we introduce GPU-accelerated tropical min-plus GEMM for shortest-path
computation on neural connectivity graphs, achieving $20.64\times$ speedup over
CPU at $N=1024$ nodes.
We evaluate across consumer integrated and discrete GPUs using the same WGSL
kernel and benchmark harness. The dated benchmark snapshot
\texttt{research/benchmark-results-2026-04-12.md} reports multi-billion
neuron-timesteps/second on discrete hardware and materially lower---but still
interactive---throughput on integrated hardware, with no shader changes across
devices. We treat those figures as dated artifact numbers and keep the
cross-hardware claim scoped to what is reproducibly logged in that snapshot.
The \cael{} hash-chain provenance adds $3.40\;\mu$s per entry median
($3.74\;\mu$s p99, 99{,}997-entry verify harness; cross-paper anchor
\texttt{research/benchmark-canonical-run.md}) across configurations.
In an embodied navigation task (harness:
\texttt{paper-snn-navigation.test.ts}, 2026-04-17), the 128-neuron agent
reaches the goal while maintaining sub-millisecond median tick time on the
harness host and producing a fully hash-chained provenance record with
100\% integrity.
Every perception--cognition--action cycle is committed to a hash-chain trace,
producing the first \emph{compute substrate} to emit auditable, tamper-evident,
replayable traces of neuromorphic decisions at hardware throughput in an
embodied simulation.
% AUDIT-2026-04-20: Scope-sharpened to disambiguate from Paper 0c (contract-framework novelty).
The \texttt{@holoscript/snn-webgpu} package is validated by a dedicated Vitest
suite in the repository (full-package run reference: 2026-04-17).
Headline GPU throughputs and tropical GEMM figures above reference the dated snapshot
\texttt{research/benchmark-results-2026-04-12.md}; interactive-scale neuron counts
(max verified discrete run: RTX~3060 Laptop GPU, \texttt{strictAdapter=true},
peak $N=1{,}048{,}576$ at $5{,}890.876$ M neurons/s;
artifact \texttt{HoloScript/.bench-logs/paper-2-lif-throughput-rtx3060.json},
generated 2026-05-06 on commit \texttt{d0a990964},
driver 591.74) are measured on
commodity integrated and discrete GPUs and run in any WebGPU-capable browser.
\end{abstract}

% ============================================================================
\section{Introduction}
\label{sec:intro}

Spiking neural networks (SNNs) offer a biologically plausible model of neural computation
in which information is encoded in the precise timing of discrete action potentials (spikes)
rather than continuous activations~\cite{maass1997networks, gerstner2002spiking}.
When coupled with physics simulations, SNNs enable \emph{embodied neuromorphic cognition}:
agents that perceive physical sensor readings, process them through biologically realistic
neural dynamics, and act upon the simulation in a closed loop.

This paradigm has been explored in robotics~\cite{tang2024braincog} and neuroscience
simulation~\cite{gewaltig2007nest}, but existing systems share a critical limitation:
they require high-performance computing clusters (NEST+Gazebo~\cite{gewaltig2007nest}),
custom neuromorphic chips (Intel Loihi~\cite{davies2018loihi}, IBM TrueNorth~\cite{merolla2014truenorth}),
or GPU-accelerated server infrastructure (BrainCog~\cite{tang2024braincog}).
No existing system delivers embodied SNN cognition in a standard web browser
with zero installation.

However, browser-native deployment is not merely a convenience feature; it is a
consequence of a deeper architectural inversion we term \emph{NN-primary,
CPU-verifier backup}. In this paradigm, spiking neural networks serve as the
primary runtime executor (Tier~1), LLM speculation serves as a warm-path
verifiable proposer (Tier~2), and CPU-side AST/typecheck/CURE verification
serves as the cold-path referee (Tier~3). The browser is the deployment
surface because WebGPU and WebNN are now production-ready across all major
engines (January 2026), making the NN-primary stack portable to commodity
hardware without datacenter dependency. On neuromorphic-class silicon this
inversion achieves single-digit milliwatt steady-state power, three orders of
magnitude below GPU-accelerated LLM inference, aligning with always-on embodied
agent workloads.

We address this gap with three contributions:

\begin{enumerate}
\item \textbf{Browser-native SNN on WebGPU as Tier~1 of an inverted runtime stack.}
We implement LIF neurons, STDP learning, multi-layer networks, and spike encoding/decoding
entirely in WGSL compute shaders, dispatched via the WebGPU API.
The system runs in Chrome, Edge, and Firefox---no plugins, no server, no compilation.
Every spike decision is committed to a \cael{} trace with cryptographic hash chaining,
producing the first tamper-evident audit trail of neuromorphic decision-making in
physics simulation and yielding a \emph{determinism receipt} that makes NN-primary
execution architecturally honest.

\item \textbf{Tropical algebra bridge between SNN and ReLU.}
We formalize the observation that LIF firing rate approximates $\relu(x)$ and show this
is not a coincidence but a consequence of the tropical max-plus semiring:
$\max(0, x) = x \tropplus 0$ where $\tropplus = \max$.
This bridge enables principled conversion between spiking and non-spiking representations.

\item \textbf{Three-tier dispatch architecture and energy analysis.}
We generalize the speculative-decoding pattern from token-stream to operation-stream,
laying out the full Tier~1 (SNN hot path), Tier~2 (LLM speculative warm path with
CPU verifier), and Tier~3 (CPU cold-path referee) architecture, and quantify the
energy gap between Tier~1 on neuromorphic-class silicon and GPU-accelerated LLM
inference.
\end{enumerate}

\paragraph{Novelty.}
This work makes four claims that, to our knowledge, have no prior instantiation
in the literature.
First, \snnwebgpu{} is the \emph{first executable Tier~1 of an inverted
runtime stack}: prior work treats SNNs as accelerators within a CPU-primary
pipeline (e.g., Loihi as a coprocessor, TrueNorth as a peripheral), whereas
our stack makes neuromorphic computation the primary executor with CPU
verification as the backup referee, generalizing the speculative-decoding
pattern from token-stream to operation-stream.
Second, \snnwebgpu{} is the \emph{first browser-native SNN runtime}: all
prior SNN frameworks require a Python ecosystem, a CUDA GPU, or a dedicated
neuromorphic chip---none run inside a standard web browser without plugins.
Third, the tropical bridge (Theorem~\ref{thm:bridge}) is the
\emph{first formal proof} that LIF rate coding approximates $\relu$ as a
direct consequence of the max-plus semiring; prior work either observed the
resemblance empirically or invoked it as a modeling heuristic without a
semiring-theoretic derivation.
Fourth, the \cael{} hash-chain provenance system produces the
\emph{first tamper-evident per-action audit trail} of neuromorphic
decision-making: every perception--cognition--action cycle is individually
hash-committed with cryptographic chaining, not merely logged at the episode
level as in all previous spiking world models and embodied SNN benchmarks.

The complete system---7 WGSL shaders, 5 TypeScript orchestration modules,
205 passing Vitest tests---is open-source as the \texttt{@holoscript/snn-webgpu} package
within the \holoscript{} ecosystem.

% ============================================================================

\paragraph{Ecosystem grounding.} This paper advances direction~D.043~(disposable neural maps, durable identity), initiative~I.002~(browser-native NPU runtime), and feedback~F.037~(papers are the product). It is a mission-load-bearing contribution to~U.002~(planet-wide prosperity via new tools).

\section{Background}
\label{sec:background}

\subsection{Leaky Integrate-and-Fire Neurons}

The LIF model~\cite{gerstner2002spiking} evolves membrane potential as $V_i(t+\Delta t) = V_\mathrm{rest} + (V_i(t)-V_\mathrm{rest})\,e^{-\Delta t/\tau} + I_\mathrm{syn}^{(i)}(t)$, where $I_\mathrm{syn}^{(j)}(t) = \sum_i W_{ji} s_i(t)$ is the weighted spike sum (equations \eqref{eq:lif}--\eqref{eq:synaptic}; full derivation in Supplemental~\S{F}).
When $V_i \geq V_\theta$, the neuron fires, resets to $V_\mathrm{reset}$, and enters a $t_\mathrm{ref}$ refractory period.
\begin{equation}
\label{eq:lif}
V_i(t + \Delta t) = V_\mathrm{rest} + \bigl(V_i(t) - V_\mathrm{rest}\bigr) \cdot e^{-\Delta t / \tau} + I_\mathrm{syn}^{(i)}(t)
\end{equation}
\begin{equation}
\label{eq:synaptic}
I_\mathrm{syn}^{(j)}(t) = \sum_{i} W_{ji} \cdot s_i(t)
\end{equation}

\subsection{WebGPU Compute Model}

WebGPU~\cite{webgpu2024spec} is a W3C standard providing low-level GPU access from web browsers via WGSL compute shaders dispatched across configurable workgroups.
Unlike WebGL, it provides read/write storage buffers, compute-only pipelines, async \texttt{mapAsync} readback, and workgroup-local shared memory---the first browser API with sufficient expressiveness for event-driven, sparse SNN simulation.

\subsection{Tropical Semirings}

A tropical semiring~\cite{maclagan2015introduction} replaces standard arithmetic: the max-plus semiring $(\mathbb{R}\cup\{-\infty\}, \max, +)$ defines $a\tropplus b=\max(a,b)$, $a\tropmult b=a+b$ (equations \eqref{eq:trop-add}--\eqref{eq:trop-mul} retained for cross-reference); the min-plus semiring uses $\min$ and encodes all-pairs shortest paths via $C_{ij}=\min_k(A_{ik}+B_{kj})$.
\begin{align}
a \tropplus b &= \max(a, b) \label{eq:trop-add} \\
a \tropmult b &= a + b \label{eq:trop-mul}
\end{align}

\subsection{CAEL Provenance System}

The Causal Agent-Environment Loop (\cael{})~\cite{holoscript2026cael} hash-chains each timestep as $E_t = H(E_{t-1}, O_t, C_t, A_t, P_t, \Delta W_t)$ using FNV-1a, so any modification to a past entry invalidates all subsequent hashes and is detectable on replay.

% ============================================================================
\section{System Architecture}
\label{sec:architecture}

\snnwebgpu{} is a browser-local perception--cognition--action loop.
Each tick encodes physics-solver sensors into rate, temporal, or delta spikes;
executes synaptic current, LIF, and optional STDP compute passes in one
\texttt{GPUCommandEncoder}; decodes output spikes into continuous actions; and
commits perception, cognition, action, and world-delta records to a \cael{}
hash chain.
The implementation uses six WGSL shader modules and 16 entry points, with
workgroup size 256 for spike/LIF kernels and $16\times16$ workgroups for dense
tropical GEMM.
The full block diagram, per-phase dataflow, shader table, and GPU resource
management notes are in the supplemental material.

% ============================================================================
\section{The Tropical Bridge: SNN Rate Coding $\approx$ ReLU}
\label{sec:tropical}

\subsection{Motivation}

A well-known empirical observation in computational neuroscience is that the
firing rate of a LIF neuron approximates the Rectified Linear Unit (ReLU)
activation function of artificial neural networks~\cite{glorot2011deep}.
Specifically, the steady-state firing rate $r(I)$ as a function of input
current $I$ is:
\begin{equation}
r(I) \approx \begin{cases}
0 & \text{if } I \leq I_\theta \\
\alpha (I - I_\theta) & \text{if } I > I_\theta
\end{cases}
\end{equation}
which is precisely $r(I) = \alpha \cdot \relu(I - I_\theta)$.
We formalize this connection through tropical algebra.

\subsection{ReLU as Tropical Addition}

In the max-plus tropical semiring $(\mathbb{R} \cup \{-\infty\}, \max, +)$,
the tropical addition operator is:
\begin{equation}
a \tropplus b = \max(a, b)
\end{equation}

The ReLU function $\relu(x) = \max(0, x)$ is therefore simply tropical addition
with the identity element:
\begin{equation}
\label{eq:relu-tropical}
\relu(x) = x \tropplus 0
\end{equation}

This is not merely a notational coincidence.
When we view neural network layers as linear algebra over the tropical semiring,
the forward pass of a ReLU network corresponds to tropical matrix multiplication:
\begin{equation}
y_j = \bigoplus_i (W_{ji} \otimes x_i) = \max_i (W_{ji} + x_i)
\end{equation}

\subsection{Formal Equivalence}

\begin{theorem}[SNN--ReLU Bridge]
\label{thm:bridge}
Let $\{s_i\}_{i=1}^N$ be spike trains from $N$ LIF neurons with steady-state
firing rates $r_i = f(I_i)$ where $f$ is the LIF f-I curve.
Let $\hat{r}_i = \frac{1}{T_w} \sum_{t=1}^{T_w} s_i(t)$ be the empirical firing
rate over time window $T_w$.
Then for sufficiently large $T_w$:
\begin{equation}
\hat{r}_i \xrightarrow{T_w \to \infty} r_i = \alpha \cdot (I_i \tropplus I_\theta) - \alpha I_\theta
\end{equation}
where $\alpha = 1/(\tau \log(V_\theta - V_\mathrm{rest}) - \tau \log(V_\mathrm{reset} - V_\mathrm{rest}))$
is the gain determined by LIF parameters.
\end{theorem}

\begin{proof}[Proof sketch]
The steady-state firing rate of a LIF neuron receiving constant current $I > I_\theta$
is $r = [\tau \log(\frac{V_\theta - V_\mathrm{rest} - I}{V_\mathrm{reset} - V_\mathrm{rest} - I})]^{-1}$.
In the suprathreshold regime where $I \gg V_\theta - V_\mathrm{rest}$,
this linearizes to $r \approx \alpha(I - I_\theta)$.
The rate coding decoder $\hat{r}_i = \text{spike\_count}/T_w$ converges to $r_i$
by the law of large numbers.
Since $\alpha(I - I_\theta)$ for $I > I_\theta$ and $0$ otherwise is exactly
$\alpha \cdot \relu(I - I_\theta)$, and $\relu(x) = x \tropplus 0$,
the equivalence holds.
\end{proof}

\subsection{Implementation: Tropical Bridge and Connectivity}

\texttt{TropicalActivationTrait} realizes Theorem~\ref{thm:bridge} in CPU and
GPU paths: the max-plus variant computes
$\alpha\max(0,x-\theta)$ and recovers $\relu{}$, while the min-plus variant
provides the cost-style dual.
The same tropical module also analyzes SNN connectivity by computing APSP via
repeated min-plus matrix multiplication,
$D=A^{\otimes\lceil\log_2 N\rceil}$ with
$(A\otimes B)_{ij}=\min_k(A_{ik}+B_{kj})$, plus a CSR SpMV path for sparse
single-source relaxations.
Full WGSL listings for \texttt{tropical\_activate},
\texttt{tropical\_min\_plus\_gemm}, and \texttt{tropical\_spmv} are in the
supplemental material.

% ============================================================================
\section{Implementation Details}
\label{sec:implementation}

\snnwebgpu{} is implemented in TypeScript (strict ES2020) with WGSL compute shaders, built with \texttt{tsup} and tested with \texttt{vitest}.
The core modules are: \texttt{gpu-context.ts} (WebGPU adapter/device lifecycle), \texttt{buffer-manager.ts} (GPU buffer allocation with staging pools), \texttt{pipeline-factory.ts} (shader module caching for 6 modules $\times$ 16 entry points), \texttt{lif-simulator.ts} (single-population LIF; 5 GPU buffers), \texttt{snn-network.ts} (multi-layer orchestrator with optional STDP), \texttt{spike-codec.ts} (rate/temporal/delta encode/decode), \texttt{TropicalActivationTrait.ts} (SNN--ReLU bridge, CPU+GPU paths), and \texttt{TropicalShortestPaths.ts} (APSP/SSSP with automatic CPU/GPU routing).

The \texttt{SNNCognitionEngine} bridges \snnwebgpu{} to the \cael{} agent loop by normalizing physical sensor readings (Pa, K, m/s$^2$) to $[0,1]$, scaling to $[0, I_\mathrm{scale}]$ mV injection current, running $\lceil \Delta t \cdot 1000 / \Delta t_\mathrm{LIF} \rceil$ LIF steps per \texttt{think()} call, and mapping aggregate firing rate to a GOAP goal stack (\textgreater30\% $\Rightarrow$ ``stabilize'', \textgreater5\% $\Rightarrow$ ``monitor'', else ``idle'').
GPU/CPU fallback to a bit-exact \texttt{CPUReferenceSimulator} is automatic when WebGPU is unavailable.
The STDP kernel uses soft-bounded Hebbian plasticity: LTP $\Delta W = \eta(1-W)$ on co-firings; LTD $\Delta W = -0.5\eta W$ when only pre fires; weights clamped to $[0,1]$.
For spike-sorted order, a wait-free Blelloch prefix-sum radix sort~\cite{blelloch1990prefix} in WGSL avoids spin-waits that deadlock on Apple/Intel GPUs.
The package is validated by 205 Vitest tests covering LIF simulation, spike codecs, SNN graph wiring, tropical kernels, GPU buffer/pipeline paths, and integration scenarios (full-package run 2026-04-17).
Full software stack details, annotated shader listings, and test coverage breakdown are in Supplemental~\S{D}.

% ============================================================================
\section{Evaluation}
\label{sec:evaluation}

\subsection{GPU Benchmark: Hardware Configurations}
\label{sec:hardware}

We evaluate the identical WGSL shader source across two adapters in one
Windows~11 hybrid notebook: Intel UHD Graphics gen-12lp integrated GPU and
NVIDIA RTX~3060 Laptop discrete GPU.
This is a cross-adapter ablation under shared CPU, DRAM, and thermals, not a
cross-site hardware study.
The full hardware table is in the supplemental material.
Benchmarks ship as \texttt{packages/snn-webgpu/benchmark-gpu.html}; a
Playwright harness writes diffable adapter/workload logs to \texttt{.bench-logs/}
for third-party reruns~\cite{snn-webgpu-harness-2026}.

\subsection{SNN Throughput}
\label{sec:lif-throughput}

The integrated baseline peaks at $1.48\times10^9$ neuron-timesteps/s for
$N=262{,}144$ before memory-bandwidth saturation, while the strict-adapter
RTX~3060 rerun reaches $5.89\times10^9$ neuron-timesteps/s at
$N=1{,}048{,}576$.
Small workloads are launch-overhead dominated; at the 1M-neuron scale the
discrete path is $4.78\times$ faster.
Both configurations clear the $>600$K neurons/s budget for 10K neurons at
60\,Hz; full per-size rows are in the supplemental material.

\subsection{Tropical GEMM Benchmarks}
\label{sec:tropical-benchmarks}

APSP uses repeated tropical GEMM squaring
($D=A^{\otimes\lceil\log_2N\rceil}$). On the Intel UHD baseline, GPU wins over
CPU from $N=128$ onward and reaches $20.64\times$ speedup at $N=1024$.
On the RTX~3060 Laptop rerun, the dense min-plus kernel reaches
$255.066$\,GFLOP/s at $N=2048$, a $27.60\times$ speedup over the integrated
GPU at the same size.
The crossover is launch-overhead dominated at small $N$ and compute-bound on
Ampere by $N\ge256$, which is exactly the SNN-connectivity workload where
online diameter/reachability analysis becomes practical during simulation.
Full CPU/GPU tables, the 2026-05-06 Dawn/Ampere run, and CSR SpMV notes are in
the supplemental material.

\subsection{CAEL Overhead}

\cael{} verify cost is CPU-bound, not GPU-bound: the canonical replay benchmark
verifies 99{,}997 entries at $3.40\,\mu$s median and $3.74\,\mu$s p99 per
entry.
At four records per agent tick, this is $13.6\,\mu$s median
($15.0\,\mu$s p99), roughly $0.09\%$ of a 60\,Hz frame.
FNV-1a is intentionally lightweight for tamper detection; SHA-256 remains an
optional stronger mode at higher cost.
The row-by-row verify table is in the supplemental material.

\iffalse
\begin{table}[t]
\centering
\caption{\cael{} FNV-1a hash-chain \emph{verify}-only throughput (\textbf{cross-paper canonical} 2026-04-18: Intel i7-11800H-class laptop, Node.js~22.x, \texttt{paper-cael-replay-benchmark.test.ts}; see \texttt{research/benchmark-canonical-run.md}).
  Median $3.40\;\mu$s/entry, p99 $3.74\;\mu$s, over 99{,}997 entries.
  \emph{Variance:} a heavier-load Vitest snapshot on 2026-04-17 showed $\sim$50\% higher per-entry medians at the same entry count (\texttt{research/benchmark-results-2026-04-17.md})---CPU scheduling, not a protocol change.
  Smaller rows use the same median for order-of-magnitude checks.
  GPU choice does not affect verify cost.}
\label{tab:cael}
\small
\measuredFrom{HoloScript/packages/engine/src/simulation/__tests__/paper-cael-replay-benchmark.test.ts}%
\measuredFrom{research/benchmark-canonical-run.md}%
\begin{tabular}{@{}rrr@{}}
\toprule
\textbf{Entries} & \textbf{Total (ms)} & \textbf{Median ($\mu$s/entry)} \\
\midrule
100 & 0.34  & 3.40 \\
1{,}000   & 3.40  & 3.40 \\
10{,}000  & 34.0  & 3.40 \\
99{,}997  & 339.6 & 3.40 \\
\bottomrule
\end{tabular}
\end{table}

Table~\ref{tab:cael} quantifies hash-chain verify cost on the canonical host.
Because FNV-1a verify is a pure scalar CPU computation, throughput depends on
the host CPU---not on which GPU serves the LIF kernel.
At the canonical operating point of 4 trace entries per agent tick
(perception, cognition, action, physics — fixed by the CAEL loop spec,
not an empirical average), the incremental verify work is exactly
$4 \times 3.40 = 13.6\;\mu$s median ($4 \times 3.74 = 15.0\;\mu$s at p99)
by arithmetic on the per-entry measurement from \texttt{paper-cael-replay-benchmark.test.ts}
(2026-04-18 canonical run, 99{,}997 entries)---about $0.082\%$ /
$0.090\%$ of the 16.6\,ms budget for 60\,Hz simulation.
The FNV-1a hash function is intentionally lightweight (not cryptographically
secure) to minimize overhead while still enabling tamper detection.
For applications requiring stronger guarantees, SHA-256 is selected via
the \texttt{useCryptographicHash} flag on \texttt{ContractConfig} at a
measured $\sim$9--24$\times$ per-hash cost on the canonical replay
harness (pure-JS FIPS~180-4 implementation; Node-native
\texttt{crypto} reduces this to $\sim$3--8$\times$ where available, but
the pure-JS path preserves universal browser compatibility).
\fi

% Twin Test Equivalence moved to supplemental (2026-05-23) for NeurIPS 9-page budget.
% Full bit-exact parity table and discussion: paper-2-snn-neurips-supplemental.tex §A.
% Summary: CPU (TypeScript) and WebGPU (WGSL) LIF implementations agree to
% <5×10⁻⁵ absolute on membrane potentials; spike masks are bit-exact across all
% tested configurations (1024/65536/2048 neurons). See Supplemental §A for
% full data~\cite{paper-0c-twin-test-2026-05-11}.

\subsection{Decoder Cost}
\label{sec:decoder-cost}

For every ``novel / faster / better'' claim, we publish the decoder (readback) cost on the same row (W.314).
The LIF kernel writes membrane potentials and spike flags to GPU storage buffers;
consuming them requires readback, but the measured cost remains below $8\%$ of
compute across the evaluated range.
At $N=65{,}536$ and 100 ticks, readback adds $8.9\,\mu$s per tick versus a
$230\,\mu$s simulation tick; the bottleneck is therefore the LIF kernel, not
the decoder.
The full row-by-row table is in the supplemental material.

\subsection{Embodied Navigation Task}
\label{sec:embodied}

In a $20\times20$ room with six obstacles, a 128-neuron SNN agent reaches the
goal at tick~38 (1.9\,s simulated time), verifies an 800-record \cael{} trace,
and runs at $0.44$\,ms mean tick time.
The baseline path is reactive and inefficient (20.0\% path efficiency), but a
supervised STDP rerun with 256 plastic synapses reaches the goal in 39 ticks
along a 29.97-unit path after 8 training episodes.
Full setup, metrics, STDP, compute-budget, and provenance tables are in the
supplemental material.

\iffalse
To demonstrate end-to-end embodied cognition, we evaluate an SNN-driven agent
in a 2D grid navigation task.

\paragraph{Setup.}
A 128 LIF neuron network navigates a $20 \times 20$ continuous room containing
6 obstacles, from start position $(2, 2)$ to goal $(18, 18)$, over 200 simulation
ticks (10.0\,s simulation time at $\Delta t = 50$\,ms per tick).
The network uses rate-coded encoding of 8 distance sensors (4 cardinal +
4 diagonal obstacle distances, plus goal bearing and goal distance) and
rate-coded decoding to 4 cardinal motor commands.
Obstacle distances are encoded as \emph{inhibitory} inputs: closer obstacles
produce stronger inhibition in the corresponding directional population.
Goal direction is encoded as \emph{excitatory} input to the aligned population.
This asymmetric encoding---inhibition from obstacles, excitation toward
goal---produces goal-directed behavior without explicit path planning.

\paragraph{Results.}

\begin{table}[t]
\centering
\caption{Navigation experiment results: 128 LIF neurons, $20 \times 20$ room,
  6 obstacles, 200 ticks. \textbf{Measured} 2026-04-17 via  \texttt{packages/engine/\allowbreak src/simulation/\allowbreak \_\_tests\_\_/\allowbreak paper-snn-navigation.test.ts}
  (CPU SNN reference backend; Node/Vitest on Windows). Mean tick 0.44\,ms ($\sim$2{,}270\,Hz); mean/median $\approx 2.7\% / 2.1\%$ of a 60\,Hz frame.}
\label{tab:nav-experiment}
\small
\measuredFrom{HoloScript/packages/engine/src/simulation/__tests__/paper-snn-navigation.test.ts}%
\begin{tabular}{@{}lr@{}}
\toprule
\textbf{Metric} & \textbf{Value} \\
\midrule
\multicolumn{2}{@{}l}{\textit{Task performance}} \\
Goal reached & Yes (tick 38, 1.9\,s sim time) \\
Final position & $(18.00, 18.00)$ \\
Distance to goal & 0.00 units \\
Path length & 160.00 units \\
Manhattan optimal & 32.00 units \\
Path efficiency & 20.0\% \\
\midrule
\multicolumn{2}{@{}l}{\textit{Neural activity}} \\
Total spikes & 37{,}312 \\
Firing rate & 14.57\% per neuron per LIF step \\
\midrule
\multicolumn{2}{@{}l}{\textit{Compute performance}} \\
Mean tick time & 0.44\,ms (median 0.34\,ms, p95 1.00\,ms) \\
Total wall time & 88\,ms (200 ticks) \\
Agent loop rate & $\sim$2{,}270\,Hz \\
Frame budget usage & median $\approx 2.1\%$; mean $\approx 2.7\%$ of 16.6\,ms (60\,Hz) \\
\midrule
\multicolumn{2}{@{}l}{\textit{CAEL provenance}} \\
Trace entries & 800 (4 per tick) \\
Trace size & 1{,}169\,KB (1{,}496 bytes/entry) \\
Hash chain integrity & 100\% valid \\
\bottomrule
\end{tabular}
\end{table}

Table~\ref{tab:nav-experiment} reports the navigation experiment results.
The SNN agent reaches the goal at tick 38 (1.9\,s simulation time) out of 200
available ticks, arriving at the target $(18, 18)$.
The traversed path length is 160.00 units versus a Manhattan optimal of 32.00,
yielding a path efficiency of 20.0\%.
This sub-optimal but successful trajectory reflects the reactive nature of the
controller: the asymmetric input encoding (obstacle inhibition + goal excitation)
produces reliable goal-directed behavior, but without memory or planning the
agent cannot anticipate obstacles and often overshoots or detours.

\paragraph{Navigation interpretation.}
Obstacle sensors inhibit blocked directions while goal bearing excites aligned
motor populations, so the baseline reaches the goal without explicit planning
but detours heavily.
A supervised STDP rerun closes that gap: a two-layer network with 256 plastic
synapses reaches the goal in $39$ ticks along a $29.97$-unit path after
8 training episodes, while the random-weight baseline fails under the same
horizon.
The same 200-tick episode emits 800 \cael{} records whose hash chain verifies.
Full STDP, compute-budget, and provenance details are in the supplemental
material.
\fi

% ============================================================================

\subsection{User Study Protocol (Preregistered)}
\label{sec:user-study}

We have prepared a pre-data protocol bundle for a within-subjects developer-diagnostic study ($N=12$ replication, pilot $N=3$) asking whether access to integrated \cael{} provenance traces and tropical-activation overlays improves a developer's ability to diagnose, verify, and trust neuromorphic-agent behaviour relative to a conventional raw-spike-raster baseline (Condition A: \snnwebgpu{} Studio viewer with CAEL trace inspection and tropical overlay; Condition B: raw spike-raster CSV + matplotlib trajectory).
Three diagnostic scenarios (obstacle overshoot, goal miss, sensor-noise drift) are drawn from the navigation task (\S\ref{sec:embodied}); order is counterbalanced with a Latin square.
Primary DV is diagnosis confidence (7-point Likert); analysis uses Wilcoxon signed-rank at $\alpha=0.05$.
This submission contributes the committed pre-data bundle only (D.011 evidence surface); it upgrades to rung~2 after IRB-approved pilot execution.
Full protocol, measures, procedure, analysis plan, and preregistration materials are in the supplemental material.
\measuredFrom{research/user-study-data/paper-2/preregistration.md}%

% ────────────────────────────────────────────────────────────────────────────
\subsection{Ablation Study}
\label{sec:ablation}

We ablate tropical activation, STDP plasticity, and \cael{} provenance against
the same WebGPU LIF benchmark harness.
The full pipeline is the only variant that jointly preserves all three
contributions: Theorem~\ref{thm:bridge}, the $39$-tick STDP navigation result,
and tamper-evident replay.
Removing tropical activation recovers only $\sim0.6\%$ throughput but breaks the
tropical consistency tests; disabling STDP returns the controller to the
reactive baseline; removing \cael{} has $<0.5\%$ throughput impact but destroys
replay integrity.
Full per-variant throughput and guarantee rows are in the supplemental material.

% ────────────────────────────────────────────────────────────────────────────
\subsection{Full-Loop Demo}
\label{sec:demo}

The shipped verifier re-runs the full navigation loop and fails unless the goal
is reached at tick~38, the trace contains 800 \cael{} records, and the hash
chain verifies.
Each provenance step commits
$h_t=\text{FNV-1a}(h_{t-1}\Vert\mathbf{s}_t\Vert\mathbf{a}_t\Vert\ell_t)$, so
tampering with a spike-rate vector, tropical activation output, or action label
invalidates replay.
The full scene instantiation, LIF dispatch, APSP, and verifier walk-through are
in the supplemental material.
\measuredFrom{research/paper-2-demo/reproduce.mjs}%
\measuredFrom{research/paper-2-demo/FULL_LOOP_VERIFICATION_2026-05-10.md}%

\section{Related Work}
\label{sec:related}

Established SNN simulation frameworks (NEST~\cite{gewaltig2007nest}, BrainCog~\cite{tang2024braincog}, snnTorch~\cite{eshraghian2023training}, SpikingJelly~\cite{fang2023spikingjelly}) require Linux clusters, CUDA GPUs, or Python ecosystems---none run in a web browser.
Browser-side neural inference via TensorFlow.js~\cite{smilkov2019tensorflow} and ONNX Runtime Web targets dense tensor operations, not spiking dynamics; the closest browser-native SNN predecessor is the WebGL-based SNN-P systems simulator of \citeauthor{cabarle2022webgl}~\cite{cabarle2022webgl}, which is limited to WebGL's compute-hostile rasterization model and does not support STDP, multi-layer LIF dynamics, or hash-chain provenance.
Neuromorphic hardware (Loihi~\cite{davies2018loihi}, TrueNorth~\cite{merolla2014truenorth}, SpiNNaker~\cite{furber2014spinnaker}) achieves superior energy efficiency but requires specialized procurement and toolchains.
The connection between tropical geometry and ReLU networks has been explored theoretically~\cite{zhang2018tropical, alfarra2022decision, maragos2021tropical}; our contribution is the first GPU-accelerated implementation that connects SNN rate codes to ReLU activations in a working system with empirical benchmarks.
Simulation provenance in scientific workflows~\cite{simmhan2005survey} and digital twins~\cite{tao2019digital} operates at the simulation level; \cael{} provides per-spike-train, per-action hash-chain provenance---a finer granularity not previously achieved.
An expanded treatment of each prior-work family is in the supplemental material.

% ============================================================================
\section{Discussion}
\label{sec:discussion}

\paragraph{Limitations.}
The LIF model omits dendritic computation, neuromodulation, and short-term plasticity; the STDP rule is a simplified Hebbian variant lacking exponential pre/post timing windows.
WebGPU availability remains incomplete on older mobile browsers, and the tropical bridge (Theorem~\ref{thm:bridge}) holds only in the suprathreshold steady-state regime.

\paragraph{Future Work.}
\label{sec:future}
STDP-learned navigation (item 1) is now closed by the measured navigation rerun in Section~\ref{sec:embodied}.
Remaining: (2) biologically accurate STDP with exponential timing windows; (3) multi-agent \cael{} traces with causal ordering; (4) Loihi~2 integration for power efficiency comparisons; (5) scaling to $10^6$ neurons via hierarchical workgroup decomposition; (6) tropical graph algorithms for SNN topology search.
Broader impact is discussed in the supplemental material.

% ============================================================================
\section{Conclusion}
\label{sec:conclusion}

We have presented \snnwebgpu{}, the first browser-native spiking neural network
system running on WebGPU compute shaders, integrated with physics simulation
via hash-chain provenance.
The system implements LIF neurons, STDP learning, 7 spike encoding/decoding
schemes, and tropical algebra operations entirely in WGSL shaders dispatched
from a standard web browser.
On the integrated baseline (Intel UHD Graphics gen-12lp), the LIF kernel
sustains $1.48 \times 10^9$ neuron-timesteps/s at $N = 262{,}144$ with
the browser-native WGSL path and no installation. The 2026-05-06 strict
adapter RTX~3060 rerun reaches $5.89 \times 10^9$ neuron-timesteps/s at
$N = 1{,}048{,}576$.
We established a formal connection between SNN rate coding and ReLU activations
through the tropical max-plus semiring, implemented as a GPU-accelerated bridge.
Our tropical min-plus GEMM achieves $20.64\times$ speedup over CPU at $N = 1024$
on integrated graphics (peak $9.88$\,GFLOP/s) and scales to $\mathbf{255.066}$
\textbf{GFLOP/s} at $N = 2048$ on the discrete GPU---a $27.60\times$ cross-GPU
speedup at the largest size---enabling online APSP analysis of SNN connectivity
graphs during simulation.
Every neural decision in the embodied agent loop is committed to a tamper-evident
hash-chain trace at a cost of $3.40\;\mu$s per entry median ($3.74\;\mu$s p99;
$\sim\!0.082\%$ / $0.090\%$ of a 60\,Hz frame budget at median/p99 when charging four trace entries per tick), producing the first compute substrate to emit auditable, replayable traces of neuromorphic
cognition at hardware throughput in physics simulation.

Browser-native SNN is deployable on integrated consumer graphics today
($1.48 \times 10^9$ neuron-timesteps/s baseline), with the same codebase,
shader, and WebGPU dispatch path scaling to a verified RTX~3060 laptop run
at $5.89 \times 10^9$ neuron-timesteps/s.
In an embodied navigation task, a 128-neuron SNN agent reaches its goal in
38 ticks (1.9\,s simulation time) at a mean tick time of 0.44\,ms---a
$\sim$2{,}270\,Hz agent loop rate with median/mean ticks $\approx 2.1\% / 2.7\%$ of a 60\,Hz frame
budget---demonstrating that goal-directed behavior emerges from asymmetric
input encoding (obstacle inhibition, goal excitation) without explicit path
planning.

The complete system---7 WGSL shaders, 16 compute entry points, 205 Vitest tests
in \texttt{@holoscript/snn-webgpu} (2026-04-17), and a multi-layer SNN orchestrator with STDP
learning---is open-source and runs in any WebGPU-capable browser with zero
installation.

More fundamentally, this paper establishes the \emph{determinism receipt} that
makes NN-primary runtime stacks architecturally honest: a spiking neural network
engine whose every decision is cryptographically traceable, reproducibly
auditable, and formally bounded by a SimulationContract verifier.
Paper~2 is therefore the credential paper for an inverted runtime in which
neuromorphic computation is the primary executor and CPU-side verification is the
backup referee---generalizing speculative decoding from token-stream to
operation-stream and opening the path to always-on embodied agents at
single-digit milliwatt power.

% ============================================================================
\section*{Reproducibility Statement}

All source code is available as the \texttt{@holoscript/snn-webgpu} package
within the HoloScript repository.
The benchmark harness (\texttt{tropical-shortest-paths.benchmark.test.ts})
is self-contained and runnable via \texttt{pnpm test}.
WGSL shader sources are provided in the supplementary materials.
The \cael{} trace format is documented in the
\texttt{CAELTrace.ts} module with full serialization/deserialization
and hash-chain verification functions.

% ============================================================================
\section*{Acknowledgments}

This work was developed within the HoloScript ecosystem by Infinitus Systems.
We thank the W3C WebGPU Community Group for the compute shader specification
that made this work possible.

% ============================================================================
% Appendix content (LIF parameters, WGSL shader listing, tropical semiring axioms,
% broader impact) is provided in the supplementary material (paper-2-snn-neurips-supplemental.tex).

{\footnotesize
\setlength{\bibsep}{0pt}
\bibliographystyle{plainnat}
\bibliography{holoscript}
}

\end{document}
